Ce que vous allez lire n'est pas un roman. Ni un polar, ni un essai,
ni une autofiction, ni un recueil de nouvelles, ni un pamphlet, ni un manuel
de savoir-vivre, ni même — précision sans doute superflue — un album pour enfants.

C'est un \textsc{ovni} — objet vittéraire non identifié.

Le texte alterne trois registres~: 1. un polar raconté à la première personne,
jalonné de nombreuses et réjouissantes scènes dialoguées~; 2. les ratiocinations du
narrateur-héros sur les sujets les plus divers~: le libre arbitre, le bel
canto, les crétins des Alpes, etc.~; et 3. un foisonnement de mentions,
références et digressions, qui forme un véritable cabinet de curiosités
historiques et culturelles. La fiction la plus \textit{gore} du récit est
dépassée par la réalité monstrueuse qui exsude du cabinet de curiosités.
L'ironie, l'humour potache et la grossière provoc — constants — ne servent
pas à enjoliver, mais sont bien la condition \textit{sine qua non} de l'\textsc{ovni}. Cerise sur le gâteau, l’auteur étouffe crânement dans l’œuf toute esquisse de début d’intrigue : «~je n’aime pas tout ce qui est intrigue, mystère et cachotterie, coup de théâtre et rebondissement de mes deux~».
